% =====================================================================
%  2.3 Andregradslikninger — ekstra utfordring
% =====================================================================

\begin{infoboks}{Velg metode selv}
Arket bygger på 2.3 Andregradslikninger. Vis regning og begrunnelse. Bruk eksakte svar.
Kontroller svarene i den opprinnelige likningen.
\end{infoboks}

\nivaahode{2}

\oppg{U1}{2}{\uh}
Noen likninger kan gjøres om til andregradslikninger. Sett $u = x^2$ og løs først likningen for $u$. Finn alle løsninger for $x$, og forklar hvor mange det er.
\begin{tasks}(2)
  \task $x^4 - 5x^2 + 4 = 0$
  \task $x^4 + x^2 - 6 = 0$
\end{tasks}

\oppg{U2}{2}{\uh}
Skriv hvilke verdier av $x$ som ikke er tillatt. Gang så med $x$ på begge sider og løs likningen.
\begin{tasks}(2)
  \task $x + \dfrac{6}{x} = 5$
  \task $2x = 3 - \dfrac{1}{x}$
\end{tasks}

\oppg{U3}{2}{\prog}
Programmet skal løse $ax^2 + bx + c = 0$.
\begin{kode}
from math import sqrt
def abc(a, b, c):
    D = b**2 - 4*a*c
    if D < 0:
        return []
    if D == 0:
        return [-b / (2*a)]
    return [(-b - sqrt(D)) / (2*a), (-b + sqrt(D)) / (2*a)]

print(abc(1, -2, -8))
print(abc(1, 2, 5))
print(abc(1, -4, 4))
\end{kode}
\begin{tasks}(1)
  \task Hva skriver programmet ut? Forklar uten å kjøre det.
  \task Hva skjer med \lstinline[style=laereverk,numbers=none,frame=none]|abc(0, 2, -4)|? Legg inn en sjekk som gjør at funksjonen gir riktig svar også da.
\end{tasks}

\nivaahode{3}

\oppg{U4}{3}{\uh}
Se på likningen $\sqrt{x + 2} = x$.
\begin{tasks}(1)
  \task Kvadrer begge sider og løs andregradslikningen du får.
  \task Sett begge løsningene inn i den opprinnelige likningen. Hva oppdager du?
  \task Forklar hvorfor en falsk løsning kan oppstå når vi kvadrerer. Bruk at $\sqrt{\phantom{x}}$ aldri er negativ.
\end{tasks}

\oppg{U5}{3}{\uh}
Lag selv en andregradslikning $ax^2 + bx + c = 0$ med løsningene $x = 2$ og $x = -5$, der $a \neq 1$.
Hvor mange slike likninger finnes? Begrunn.

\oppg{U6}{3}{\uh}
En andregradslikning $ax^2 + bx + c = 0$ har løsningene $x_1$ og $x_2$.
\begin{tasks}(1)
  \task Bruk abc-formelen til å vise at $x_1 + x_2 = -\dfrac{b}{a}$ og $x_1 \cdot x_2 = \dfrac{c}{a}$.
  \task Bruk dette til å forklare hvorfor sum--produkt-metoden virker for $x^2 + bx + c = 0$.
\end{tasks}

\begin{oppsummering}
Se på likningen først: felles faktor, kvadratrot eller kvadratsetning er ofte raskere enn abc-formelen.
Likninger kan gjøres om til andregradslikninger, men sjekk alltid løsningene i den opprinnelige likningen.
Løsningene har fast sum $-\frac{b}{a}$ og fast produkt $\frac{c}{a}$.
\end{oppsummering}
